\documentclass{article}
\usepackage{amsmath,hyperref}
\title{Certified time bounds for wider protected corridor intervals}
\date{}
\begin{document}
\maketitle
Keep the ordered speed-only corridor and throughout coverage from
\href{https://github.com/eternal-november/eos-research-programmes/pull/409}{PR~409}.
For an ascending leg $A\le B$, let $S$ be the set of placements
satisfying the fixed coordinate bounds, clearance, robust radio
connectivity and pointwise robust sensing of the closed interval.
In particular, $S$ contains both endpoints. It is a finite union of
closed difference-constraint cells indexed by sensing end witnesses,
and is closed under coordinatewise max. The known min/max
decomposition handles noncomparable endpoints separately.
No new physical disturbance is allowed during a planned motion.

\paragraph{Greatest next sample.}
Given an upper vector $u$, compute the greatest point of
$S\cap[A,u]$, or report it empty. In each cell initialize the
upper vector with $u$ and a zero reference coordinate. Each inequality
$x_j-x_i\le c$ lowers the current upper bound of $x_j$ to at most
the current upper bound of $x_i$ plus $c$. The lower bounds $A$
are also inequalities using the reference coordinate. Every feasible
point stays below the current bounds. A negative cycle or a decrease
below zero at the reference proves that the cell is empty.
Otherwise ordinary shortest-path propagation gives a fixed point
that satisfies all constraints and dominates every feasible point
in the cell. Take the coordinatewise maximum of the cell optima.
It is still feasible by closure, so it is the greatest point
of the union under the same cap. Each of the quadratically many candidate cells has a linear-size
description, so a projection takes polynomial work.

For a mesh $h>0$, from $x_0=A$ set
\[
 x_{t+1}=\max\bigl(S\cap[A,\min(B,x_t+h v)]\bigr).
\]
The old sample is admissible, so this step is defined and never moves
any coordinate backward. The greatest sample dominates the sample
at the same time of every physically feasible ascending trajectory:
induct on the sample number, using the speed cap and the fact that
the actual trajectory's samples lie in $S$. If the rollout first
reaches $B$ at step $m>0$, every real motion takes more than
$(m-1)h$. If it has not arrived by step $N$, no real motion meets
deadline $Nh$. If the iteration reaches a fixed point other than $B$,
it rules out every finite arrival time, not only that mesh horizon.

\paragraph{A positive certificate requires more than safe samples.}
Consecutive samples have speed-feasible affine chords, but their
sensing can go unsafe between samples. A chord is covered exactly
when every survivor set still covers the protected interval at all
endpoint-order events and one point in every open event stratum.
The sensing interval endpoints are affine during a chord, so their
relative order and their order against $L,U$ do not change within
a stratum. Clearance and radio bounds interpolate directly. Where
all chords pass, $mh$ is a certified upper bound; the gap to the
lower bound is at most $h$. Else the result is lower-bound-only.
On the disconnected clearance-5 example from PR~409, the rollout
with step 6 appears to arrive across a gap; the chord check rejects
that motion. At step $1/4$, even sample arrival fails. This avoids
replacing an invariant by a sample test.

As meshes tend to zero, lower bounds from the unrestricted
sample rollout converge to the true optimum, with infinity allowed.
If a sequence of meshes could reach strictly earlier by a fixed
margin, the padded affine sample interpolants would be uniformly
bounded and uniformly Lipschitz on that shorter horizon. Their
distance to $S$ is bounded by a constant times the mesh. A uniformly
convergent subsequence would give a real path in closed $S$ within
that horizon, a contradiction. This argument does not give a
convergence rate, and it does not say the chords eventually pass.
There is no finite complete minimum-time algorithm asserted for
the general width here.

\paragraph{A broad zone invalidates the old crossing order.}
For ordinary single coverage of $[5,9]$ with radius 5, gap 1 and
radio range 10, move from $(-2,-1,7)$ to $(5,6,7)$; the third
robot stays at 7 and covers the whole zone. Unit speeds on the
first two make 7 optimal by translation. The narrow-interval rule
would require robot 2 to cross 4 before robot 1 crosses 0.
That spurious precedence forces at least 10 and attains 10.
In another tested fixture with four robots, double coverage of
$[0,9]$ and gap 3, the same erroneous precedence rules out a
feasible 3.5-unit trajectory. The exact full cell-route comparator
checks both. The width assumption matters even when endpoints,
radio and clearances all agree.

The matched reference enumerates routes in the same convex cells,
using the general \href{https://arxiv.org/abs/2205.04422}{graphs of
convex sets} representation; its numerical LP output is a replay
comparison, not the proof. No runtime superiority, field validation
or R6/I2 closure is claimed.
\end{document}
